%%%%%%%%%%%%%%%%%%%%%%%%%%%%%%%%%%%%%%%%%%%%%%%%%%%%%%%%%%%%%%%
% Plantilla genérica para documento académico / TFM / tesis
%%%%%%%%%%%%%%%%%%%%%%%%%%%%%%%%%%%%%%%%%%%%%%%%%%%%%%%%%%%%%%%

\documentclass[12pt]{article}

% Paquetes principales
\usepackage[spanish, es-tabla]{babel}
\usepackage{titlesec}
\usepackage{setspace}
\usepackage{graphicx}
\usepackage[hidelinks]{hyperref}
\usepackage{array}
\usepackage{times}
\usepackage{upgreek}
\usepackage{todonotes}
\usepackage{url}
\usepackage{longtable}
\usepackage{pdflscape}

% Ruta de imágenes
\graphicspath{{figures/}{imgs/}}

% Nombre del índice
\renewcommand{\contentsname}{Índice de contenido}

% Espaciado
\doublespacing

\begin{document}

%%%%%%%%%%%%%%%%%%%%%%%%%%%%%%%%%%%%%%%%%%%%%%%%%%%%%%%%%%%%%%%
% Portada
%%%%%%%%%%%%%%%%%%%%%%%%%%%%%%%%%%%%%%%%%%%%%%%%%%%%%%%%%%%%%%%

\begin{titlepage}
    \begin{center}
        \vspace*{1cm}

        \Huge
        \textbf{Título del trabajo académico}\\

        \vspace{0.5cm}

        \LARGE
        Subtítulo del trabajo, si aplica

        \vspace{3cm}

        \textbf{Alumno/a:}\\
        Nombre del estudiante\\

        \vspace{0.5cm}

        \textbf{Tutor/a:}\\
        Nombre del tutor o tutora\\

        \vspace{1cm}

        \textbf{Institución:}\\
        Nombre de la universidad o institución\\

        \vspace{0.5cm}

        \textbf{Programa académico:}\\
        Nombre del programa, máster, grado o asignatura\\

        \vfill

        \Large
        Lugar y fecha

    \end{center}
\end{titlepage}

%%%%%%%%%%%%%%%%%%%%%%%%%%%%%%%%%%%%%%%%%%%%%%%%%%%%%%%%%%%%%%%
% Páginas preliminares
%%%%%%%%%%%%%%%%%%%%%%%%%%%%%%%%%%%%%%%%%%%%%%%%%%%%%%%%%%%%%%%

\pagenumbering{roman}

\begin{center}
\LARGE
\textbf{Agradecimientos}
\end{center}

Escriba aquí los agradecimientos del trabajo.

\newpage

\begin{center}
\LARGE
\textbf{Resumen}
\end{center}

Escriba aquí el resumen del trabajo. Debe presentar brevemente el problema abordado, el objetivo principal, la metodología utilizada, los resultados más relevantes y las conclusiones generales.

\vspace{0.5cm}

\textbf{Palabras clave:} palabra clave 1, palabra clave 2, palabra clave 3, palabra clave 4.

\newpage

%%%%%%%%%%%%%%%%%%%%%%%%%%%%%%%%%%%%%%%%%%%%%%%%%%%%%%%%%%%%%%%
% Índices
%%%%%%%%%%%%%%%%%%%%%%%%%%%%%%%%%%%%%%%%%%%%%%%%%%%%%%%%%%%%%%%

\tableofcontents
\newpage

\listoffigures
\listoftables
\newpage

%%%%%%%%%%%%%%%%%%%%%%%%%%%%%%%%%%%%%%%%%%%%%%%%%%%%%%%%%%%%%%%
% Cuerpo del documento
%%%%%%%%%%%%%%%%%%%%%%%%%%%%%%%%%%%%%%%%%%%%%%%%%%%%%%%%%%%%%%%

\pagenumbering{arabic}

\section{Introducción}

En esta sección se presenta el contexto general del tema de estudio, la problemática abordada, la importancia del trabajo y una visión general de la estructura del documento.

\subsection{Justificación y antecedentes}

Explique aquí la relevancia del tema, los antecedentes principales, el problema identificado y la necesidad de desarrollar el estudio.

\begin{figure}[h]
\centering
\includegraphics[width=1.0\textwidth]{nombre-imagen.jpg}
\caption{Descripción de la figura}
\label{fig:ejemplo}
\end{figure}

\subsection{Objetivos del estudio}

\subsubsection{Objetivo general}

Describir el objetivo general del trabajo.

\subsubsection{Objetivos específicos}

\begin{itemize}
    \item Describir el primer objetivo específico.
    \item Describir el segundo objetivo específico.
    \item Describir el tercer objetivo específico.
\end{itemize}

\newpage

%%%%%%%%%%%%%%%%%%%%%%%%%%%%%%%%%%%%%%%%%%%%%%%%%%%%%%%%%%%%%%%
% Materiales y métodos
%%%%%%%%%%%%%%%%%%%%%%%%%%%%%%%%%%%%%%%%%%%%%%%%%%%%%%%%%%%%%%%

\section{Materiales y Métodos}

En esta sección se describen los materiales, herramientas, tecnologías, fuentes de información y métodos utilizados para desarrollar el trabajo.

\subsection{Estado del arte}

Presente aquí una revisión de trabajos previos, enfoques existentes, herramientas relacionadas y fundamentos teóricos relevantes.

\subsubsection{Tema o enfoque relacionado}

Describa aquí un tema específico del estado del arte.

\subsection{Metodología}

Explique el procedimiento seguido para desarrollar el trabajo, incluyendo fases, técnicas, criterios de análisis y procesos de validación.

\subsubsection{Descripción del método}

Detalle el método utilizado, sus etapas y la forma en que se aplicó en el desarrollo del proyecto.

\subsubsection{Caracterización de datos, repositorios o fuentes de información}

Describa aquí cómo se analizaron, clasificaron o caracterizaron los datos, documentos, repositorios, registros u otros elementos utilizados.

\subsubsection{Desarrollo e implementación de la herramienta o solución}

Explique aquí el diseño, desarrollo, arquitectura, integración y funcionamiento de la herramienta, sistema, modelo o propuesta desarrollada.

\newpage

%%%%%%%%%%%%%%%%%%%%%%%%%%%%%%%%%%%%%%%%%%%%%%%%%%%%%%%%%%%%%%%
% Resultados
%%%%%%%%%%%%%%%%%%%%%%%%%%%%%%%%%%%%%%%%%%%%%%%%%%%%%%%%%%%%%%%

\section{Resultados}

En esta sección se presentan los resultados obtenidos durante el desarrollo, implementación o evaluación del trabajo.

\subsection{Planificación del proyecto}

Describa la planificación realizada, fases del proyecto, cronograma, tareas principales o metodología de gestión utilizada.

\subsection{Análisis de datos, repositorios o casos de estudio}

Presente los resultados del análisis realizado sobre los datos, repositorios, documentos o casos evaluados.

\subsection{Ejemplo de uso de la solución desarrollada}

Incluya un ejemplo práctico que muestre cómo se utiliza la herramienta, metodología o propuesta desarrollada.

\subsection{Evaluación de los resultados}

Presente los criterios de evaluación, métricas utilizadas, resultados obtenidos y análisis de la efectividad de la solución propuesta.

\newpage

%%%%%%%%%%%%%%%%%%%%%%%%%%%%%%%%%%%%%%%%%%%%%%%%%%%%%%%%%%%%%%%
% Conclusiones y trabajo futuro
%%%%%%%%%%%%%%%%%%%%%%%%%%%%%%%%%%%%%%%%%%%%%%%%%%%%%%%%%%%%%%%

\section{Conclusiones y Trabajo Futuro}

En esta sección se resumen los principales hallazgos del trabajo, el cumplimiento de los objetivos y las aportaciones realizadas.

\subsection{Discusión y limitaciones}

Analice los resultados obtenidos, sus implicaciones, posibles limitaciones metodológicas, técnicas o prácticas, y aspectos que deben considerarse al interpretar los hallazgos.

\subsection{Trabajo futuro}

Describa posibles mejoras, nuevas líneas de investigación, funcionalidades adicionales o próximos pasos derivados del trabajo realizado.

\newpage

%%%%%%%%%%%%%%%%%%%%%%%%%%%%%%%%%%%%%%%%%%%%%%%%%%%%%%%%%%%%%%%
% Bibliografía
%%%%%%%%%%%%%%%%%%%%%%%%%%%%%%%%%%%%%%%%%%%%%%%%%%%%%%%%%%%%%%%

\bibliographystyle{plain}
\bibliography{export}

\end{document}